% ch43_quasi_categories.tex — Volume IV Chapter 7

\chapter{Quasi-Categories: Definition and First Examples}

\begin{overviewbox}
A \term{quasi-category} is a simplicial set in which every inner horn
admits a filler. The condition is one weakening of the nerve
characterization (Theorem~39.\ref{thm:nerve-characterization}) — drop
the uniqueness; keep only existence — and the resulting structure is
the working model of an $\infty$-category that organizes Volume~IV's
remaining chapters.

This chapter introduces the definition, builds the dictionary between
quasi-categorical and ordinary-categorical vocabulary, and works through
the first examples. We meet:
\begin{itemize}
  \item Nerves of ordinary categories (Chapter~39): quasi-categories with
        unique fillers.
  \item Kan complexes (Chapter~39): quasi-categories whose mapping spaces
        are all $\infty$-groupoids.
  \item Differential graded categories (locally) and topological
        categories (after a coherent-nerve construction): the bridge to
        practice.
\end{itemize}

The technical heart of the chapter is the construction of the
\term{homotopy category} $\mathrm{ho}(\mathcal{C})$ of a quasi-category
$\mathcal{C}$, an ordinary category whose objects are the $0$-simplices of
$\mathcal{C}$ and whose morphisms are equivalence classes of $1$-simplices
under a specific homotopy relation defined via $2$-simplex fillers. The
adjunction $\tau_1 \dashv \mathrm{N}$ of Chapter~39 restricts to a
functor $\mathrm{ho} : \mathbf{QCat} \to \mathbf{Cat}$ — the first
descent from the $\infty$-categorical to the $1$-categorical world.

Mapping spaces, equivalences, the Joyal model structure, joins/slices,
and the fibrations of Cluster~D all build on the basic combinatorics
established here.
\end{overviewbox}


% ═══════════════════════════════════════════════════════════════════════════
\section{Quasi-Categories}
\label{sec:quasi-categories}
% ═══════════════════════════════════════════════════════════════════════════

\subsection{Informally}

A quasi-category is what one gets when ``a category up to coherent
homotopy'' is presented combinatorially. Take a simplicial set and read
its $1$-simplices as morphisms, its $0$-simplices as objects. The
\emph{composition} of two composable $1$-simplices $f, g$ is not a single
$1$-simplex but a $2$-simplex whose two short edges are $f, g$ and whose
long edge is the composite. Such a $2$-simplex always exists (this is
exactly the inner-horn filling condition), but it is not unique — different
fillers give different composites, all related by higher-dimensional
simplices.

\subsection{Formalizing}

\begin{definition}[Quasi-category]
\label{def:quasi-category}
A \term{quasi-category} (also called an \term{$\infty$-category} or, in
older literature, \term{weak Kan complex}) is a simplicial set
$\mathcal{C}$ in which every inner horn admits a filler: for $n \geq 2$
and $0 < k < n$, every map $\Lambda^n_k \to \mathcal{C}$ extends to
$\Delta^n \to \mathcal{C}$.
\end{definition}

\begin{howtosay}
$\mathcal{C}$ \sayit{``a quasi-category $\mathcal{C}$''} or
\sayit{``an $\infty$-category $\mathcal{C}$''}; we use both interchangeably,
preferring ``quasi-category'' when emphasizing the simplicial-set
combinatorics and ``$\infty$-category'' when emphasizing the categorical
content. The category $\mathbf{QCat}$ is the full subcategory of
$\mathbf{sSet}$ on quasi-categories.
\end{howtosay}

\begin{example_thm}[Nerves of categories]
\label{ex:nerves-quasi}
Every nerve $\mathrm{N}(\mathcal{D})$ of an ordinary category $\mathcal{D}$
is a quasi-category. The inner-horn filler is in fact unique
(Theorem~39.\ref{thm:nerve-characterization}), so nerves are
``$1$-truncated'' quasi-categories. Conversely, a quasi-category with
unique inner-horn fillers is a nerve.
\end{example_thm}

\begin{example_thm}[Kan complexes]
\label{ex:kan-quasi}
Every Kan complex (Definition~39.\ref{def:kan-complex}) is a
quasi-category: inner horns are a subset of all horns. In particular,
$\mathrm{Sing}(T)$ for any topological space $T$ is a quasi-category —
specifically, an $\infty$-groupoid (every morphism is invertible up to
coherent homotopy).
\end{example_thm}

\begin{example_thm}[Joyal model fibrant simplicial sets]
\label{ex:joyal-fibrant}
The fibrant objects of the Joyal model structure on $\mathbf{sSet}$
(Remark~41.\ref{rem:joyal-preview}) are exactly the quasi-categories.
This is the precise sense in which $\mathbf{sSet}$ ``presents''
$\infty$-categories: $\mathrm{Ho}(\mathbf{sSet}_{\mathrm{Joyal}})$ is
the homotopy category of $\infty$-categories.
\end{example_thm}

\begin{remark}[Why drop uniqueness?]
Ordinary category theory enforces that composites are unique by axiom.
In ``higher-dimensional'' settings — chain complexes up to chain homotopy,
spaces up to coherent homotopy, derived categories with their triangulated
structure — composites are unique only up to a higher-dimensional
equivalence, which itself is unique up to a still-higher equivalence,
and so on. Quasi-categories present this entire tower with a single
combinatorial structure: the simplices of dimension $\geq 2$ are the
witnesses to all the higher coherences.
\end{remark}


% ═══════════════════════════════════════════════════════════════════════════
\section{The Homotopy Category}
\label{sec:homotopy-category-qcat}
% ═══════════════════════════════════════════════════════════════════════════

\subsection{Formalizing}

A quasi-category has too much higher-dimensional data to be a category in
the ordinary sense — but it has a quotient that is.

\begin{definition}[Homotopy between $1$-simplices]
\label{def:homotopy-1simplex}
Let $\mathcal{C}$ be a quasi-category and $f, g : x \to y$ two
$1$-simplices with the same source $x$ and target $y$ (i.e., $d_1 f =
d_1 g = x$, $d_0 f = d_0 g = y$). We say $f$ is \term{homotopic} to $g$,
written $f \sim g$, if there exists a $2$-simplex $\sigma : \Delta^2 \to
\mathcal{C}$ whose boundary is
\begin{equation*}
  d_0 \sigma = g, \quad d_1 \sigma = f, \quad d_2 \sigma = s_0 x,
\end{equation*}
i.e., a $2$-simplex with edges $g$ (face opposite $0$), $f$ (face opposite
$1$), and the degenerate identity $s_0 x$ on $x$ (face opposite $2$).
\end{definition}

\begin{lemma}[Homotopy is an equivalence relation]
\label{lem:homotopy-eq-rel}
The relation $\sim$ on $1$-simplices with fixed source and target is an
equivalence relation: reflexive, symmetric, and transitive.
\end{lemma}

\begin{proof}[Proof sketch]
\emph{Reflexive.} Use the degenerate $s_1 f : \Delta^2 \to \mathcal{C}$
with $d_0(s_1 f) = f$, $d_1(s_1 f) = f$, $d_2(s_1 f) = s_0 x$.

\emph{Symmetric.} Given a witnessing $2$-simplex for $f \sim g$ together
with the degenerate $s_1 g$ for $g \sim g$, fill the inner horn
$\Lambda^3_2$ they form to produce a witness for $g \sim f$.

\emph{Transitive.} Given witnesses for $f \sim g$ and $g \sim h$,
together with $s_0 x$, fill the inner horn $\Lambda^3_1$ to produce a
witness for $f \sim h$. \qed
\end{proof}

\begin{definition}[Homotopy category]
\label{def:ho-cat}
The \term{homotopy category} $\mathrm{ho}(\mathcal{C})$ of a
quasi-category $\mathcal{C}$ is the ordinary category with:
\begin{itemize}
  \item Objects: $\mathrm{Ob}(\mathrm{ho}(\mathcal{C})) = \mathcal{C}_0$.
  \item Morphisms $\mathrm{ho}(\mathcal{C})(x, y) = \mathcal{C}_1(x, y) / \sim$,
        equivalence classes of $1$-simplices with source $x$, target $y$.
  \item Composition: for $[f] : x \to y$ and $[g] : y \to z$, choose
        $2$-simplex fillers of the inner horn $\Lambda^2_1 \to \mathcal{C}$
        on $(f, g)$ and set $[g] \circ [f] = [d_1 \sigma]$ where $\sigma$
        is any such filler.
  \item Identity: $[s_0 x] : x \to x$.
\end{itemize}
\end{definition}

\begin{theorem}[Well-definedness of $\mathrm{ho}$]
\label{thm:ho-welldef}
Composition in $\mathrm{ho}(\mathcal{C})$ does not depend on the choice
of $2$-simplex filler, nor on the representatives $f, g$ within
equivalence classes. Associativity and unit laws hold strictly in
$\mathrm{ho}(\mathcal{C})$.
\end{theorem}

\begin{proof}[Proof sketch]
\emph{Filler-independence.} Two fillers $\sigma, \sigma'$ of the same
inner horn glue into an inner horn $\Lambda^3_1$ whose filling witnesses
$d_1 \sigma \sim d_1 \sigma'$.

\emph{Representative-independence.} Replacing $f$ by $f'$ with $f \sim f'$
substitutes the relevant $2$-simplex faces by homotopic ones; the
$\Lambda^3$ horn-filling argument transports the conclusion.

\emph{Associativity.} Given three composable $1$-simplices, fill a
$\Lambda^4_2$ horn with the relevant $2$-cells as faces. The $4$-simplex
filler witnesses associativity.

\emph{Units.} Degenerate $1$-simplices $s_0 x$ are units for the
composition by direct simplex-counting.
\qed
\end{proof}

\begin{corollary}[Functor $\mathrm{ho} : \mathbf{QCat} \to \mathbf{Cat}$]
\label{cor:ho-functor}
The assignment $\mathcal{C} \mapsto \mathrm{ho}(\mathcal{C})$ extends to a
functor $\mathrm{ho} : \mathbf{QCat} \to \mathbf{Cat}$, left adjoint to
the nerve $\mathrm{N} : \mathbf{Cat} \to \mathbf{QCat}$.
\end{corollary}

\begin{proof}[Proof sketch]
Extend $\tau_1 \dashv \mathrm{N}$ of Theorem~39.\ref{thm:nerve-tau1-adj}
to the full subcategory $\mathbf{QCat}$; the adjunction restricts.
$\mathrm{ho}(\mathcal{C}) = \tau_1(\mathcal{C})$ for $\mathcal{C}$ a
quasi-category. \qed
\end{proof}

\begin{remark}[``Homotopy $1$-category'']
The functor $\mathrm{ho}$ flattens an $\infty$-category to its underlying
$1$-category by collapsing all higher coherence data. The kernel of this
flattening is exactly the higher-categorical content that
distinguishes quasi-categories from ordinary categories.
\end{remark}


% ═══════════════════════════════════════════════════════════════════════════
\section{Bridges and Working Examples}
\label{sec:qcat-bridges}
% ═══════════════════════════════════════════════════════════════════════════

\subsection{Differential graded categories}

\begin{remark}[The dg-nerve]
\label{rem:dg-nerve}
For a differential graded (dg) category $\mathcal{D}$ over a ring $k$, the
\term{dg-nerve} $\mathrm{N}^{\mathrm{dg}}(\mathcal{D})$ is a quasi-category
whose $n$-simplices are coherent diagrams in $\mathcal{D}$ together with
$\mathbb{E}_\infty$-multiplicative homotopies between their composites.
This construction sends dg-categories to quasi-categories in a way
compatible with the homotopy theory of dg-categories. It is the standard
source of \emph{stable} quasi-categories arising in homological algebra
(Chapter~53).
\end{remark}

\subsection{Topological categories}

\begin{remark}[The coherent nerve]
\label{rem:coherent-nerve}
A \term{topological category} (or \term{simplicial category}) is a
category enriched in $\mathbf{Top}$ (or $\mathbf{sSet}$). Lurie's
\term{coherent nerve}
$\mathrm{N}^{\mathrm{coh}} : \mathbf{sCat} \to \mathbf{sSet}$ sends each
simplicial category to a simplicial set whose $n$-simplices are
``coherent'' chains of morphisms. The image of the coherent nerve lands
in quasi-categories, and the functor descends to an equivalence of
$\infty$-categories
\begin{equation*}
  \mathbf{sCat} \;\simeq_{\mathrm{Q}}\; \mathbf{sSet}_{\mathrm{Joyal}}
\end{equation*}
(a Quillen equivalence) — saying that the two definitions of
``$\infty$-category'' (simplicial categories vs.\ quasi-categories)
agree.
\end{remark}

\subsection{The $\infty$-category of spaces}

\begin{example_thm}[The $\infty$-category $\mathcal{S}$ of spaces (anima)]
\label{ex:infty-spaces}
The \term{$\infty$-category of spaces} (or \term{anima}) is the
quasi-category
\begin{equation*}
  \mathcal{S} \;=\; \mathrm{N}^{\mathrm{coh}}(\mathbf{Kan}),
\end{equation*}
the coherent nerve of the simplicial category of Kan complexes (with
mapping spaces $\mathrm{Map}_{\mathbf{sSet}}(X, Y)$). Equivalently,
$\mathcal{S}$ is the $\infty$-categorical localization of $\mathbf{sSet}$
at weak equivalences. Its homotopy category $\mathrm{ho}(\mathcal{S})
\simeq \mathrm{Ho}(\mathbf{Top})$ is the classical homotopy category of
spaces.
\end{example_thm}

\begin{remark}[Notation for spaces in $\mathcal{S}$]
We will write objects of $\mathcal{S}$ as \emph{anima} (Lurie's term),
\emph{spaces}, or simply ``$\infty$-groupoids,'' depending on emphasis.
The notation $\mathcal{S}$ for the $\infty$-category of spaces is
standard in Lurie's \emph{Higher Topos Theory} and \emph{Kerodon}.
\end{remark}

\subsection{Three pictures of an $\infty$-category}

\begin{remark}[Three models, one homotopy theory]
There are at least three established combinatorial models for
$\infty$-categories:
\begin{itemize}
  \item \emph{Quasi-categories} (Joyal–Lurie): simplicial sets with
        inner-horn fillers. Our working model in Volume~IV.
  \item \emph{Simplicial categories} ($\mathbf{sCat}$): categories
        enriched in $\mathbf{sSet}$.
  \item \emph{Complete Segal spaces} (Rezk): bisimplicial sets satisfying
        completeness conditions.
\end{itemize}
All three are connected by Quillen equivalences (Bergner, Joyal–Tierney,
Lurie); the choice of model is a matter of which combinatorics is most
convenient. Volume~IV exclusively uses quasi-categories; readers
graduating to Lurie's \emph{Higher Topos Theory} encounter all three.
\end{remark}


% ═══════════════════════════════════════════════════════════════════════════
\section{Exercises}
\label{sec:qcat-exercises}
% ═══════════════════════════════════════════════════════════════════════════

\begin{exercisesbox}
\begin{qa}
\qaitem{Define a quasi-category.}{A simplicial set in which every inner horn $\Lambda^n_k \to \mathcal{C}$ (with $0 < k < n$) extends to $\Delta^n \to \mathcal{C}$.}
\qaitem{What's the relation to nerves?}{Nerves of ordinary categories are exactly the quasi-categories with \emph{unique} inner-horn fillers. Dropping uniqueness gives quasi-categories proper.}
\qaitem{Why are Kan complexes quasi-categories?}{Inner horns are a subset of all horns; if every horn fills, in particular every inner one does.}
\qaitem{What's the heuristic for inner-horn filling?}{Composition: $\Lambda^2_1 \to \mathcal{C}$ is two composable $1$-simplices; the filler $\Delta^2 \to \mathcal{C}$ provides a composite as the long edge, plus a $2$-cell witnessing it.}
\qaitem{When is composition unique?}{Never strictly — only up to a $\Lambda^3_1$-filler-witnessed homotopy. In the nerve of an ordinary category, this homotopy is forced to be the identity by the uniqueness of fillers.}
\qaitem{Define homotopy of $1$-simplices.}{$f \sim g$ iff there is a $\Delta^2 \to \mathcal{C}$ with edges $f$, $g$, and $s_0 x$ on $x = d_1 f$.}
\qaitem{Why is $\sim$ an equivalence relation?}{Reflexivity uses degenerate $2$-simplices; symmetry and transitivity use $\Lambda^3$-filler arguments built from the given witnesses plus degeneracies.}
\qaitem{Construct $\mathrm{ho}(\mathcal{C})$.}{Objects $\mathcal{C}_0$, morphisms $\mathcal{C}_1/\sim$, composition by inner-horn filler (well-defined up to $\sim$), identities $[s_0 x]$.}
\qaitem{Why is composition in $\mathrm{ho}(\mathcal{C})$ well-defined?}{Two fillers of the same horn produce homotopic long edges (by a $\Lambda^3_1$ filler); replacing $f, g$ by homotopic representatives also produces homotopic composites.}
\qaitem{Why does associativity hold in $\mathrm{ho}(\mathcal{C})$?}{Three composable $1$-simplices plus their two composites give the data for a $\Lambda^4_2$-horn, whose filler is a $4$-simplex witnessing associativity at the homotopy level.}
\qaitem{State the adjunction $\mathrm{ho} \dashv \mathrm{N}$.}{$\mathbf{Cat}(\mathrm{ho}(\mathcal{C}), \mathcal{D}) \cong \mathbf{QCat}(\mathcal{C}, \mathrm{N}(\mathcal{D}))$. Maps into a $1$-category factor through the homotopy $1$-category.}
\qaitem{What does $\mathrm{ho}$ destroy?}{All higher-dimensional simplices beyond their incidence with $1$-simplices — i.e., all higher coherence data.}
\qaitem{Name three models of $\infty$-categories.}{Quasi-categories (Joyal–Lurie); simplicial categories ($\mathbf{sSet}$-enriched); complete Segal spaces (Rezk).}
\qaitem{Why does Volume~IV use quasi-categories?}{They are the most combinatorially direct: simplicial sets with a single horn-filling condition. The other models are Quillen-equivalent but require more setup.}
\qaitem{What is the dg-nerve?}{$\mathrm{N}^{\mathrm{dg}}(\mathcal{D})$: a quasi-category from a dg-category, whose simplices encode coherent chain homotopies. The source of stable $\infty$-categories in homological algebra.}
\qaitem{What is the coherent nerve?}{$\mathrm{N}^{\mathrm{coh}} : \mathbf{sCat} \to \mathbf{sSet}$ sending a simplicial category to a quasi-category. Lands in $\mathbf{QCat}$; descends to a Quillen equivalence $\mathbf{sCat} \simeq_{\mathrm{Q}} \mathbf{sSet}_{\mathrm{Joyal}}$.}
\qaitem{What is $\mathcal{S}$?}{The $\infty$-category of spaces (anima); $\mathcal{S} = \mathrm{N}^{\mathrm{coh}}(\mathbf{Kan})$. Equivalently, the $\infty$-localization of $\mathbf{sSet}$ at weak equivalences.}
\qaitem{What's $\mathrm{ho}(\mathcal{S})$?}{The classical homotopy category $\mathrm{Ho}(\mathbf{Top})$ of topological spaces (or equivalently, $\mathrm{Ho}(\mathbf{sSet})$).}
\qaitem{Distinguish ``$\infty$-category'' and ``quasi-category.''}{Same thing in this book. ``Quasi-category'' emphasizes simplicial-set combinatorics; ``$\infty$-category'' emphasizes the categorical content.}
\qaitem{What is an ``$\infty$-groupoid''?}{A quasi-category in which every morphism is an equivalence — equivalently, a Kan complex. The objects of $\mathcal{S}$.}
\qaitem{Why doesn't quasi-category composition simply pick a chosen filler?}{Functoriality would require coherent choice across all horns, and no such global choice exists in general — choosing strictly would specialize to nerves of categories.}
\qaitem{What's the next chapter's contribution?}{Mapping spaces $\mathrm{Map}_\mathcal{C}(x, y)$ as Kan complexes; equivalences of quasi-categories; the Joyal model structure that organizes them.}
\qaitem{What's the philosophical takeaway?}{The same simplicial set $\mathbf{sSet}$ supports two model structures: Kan–Quillen (for spaces) and Joyal (for $\infty$-categories). The combinatorics is shared; the homotopy theory differs in which class of morphisms is inverted.}
\end{qa}
\end{exercisesbox}


% ═══════════════════════════════════════════════════════════════════════════
\section*{Chapter Summary}
\addcontentsline{toc}{section}{Chapter Summary}
% ═══════════════════════════════════════════════════════════════════════════

\begin{summarybox}
\textbf{Quasi-category.}\quad A simplicial set in which every inner horn
$\Lambda^n_k \to \mathcal{C}$ ($0 < k < n$) extends to $\Delta^n \to
\mathcal{C}$. Nerves of categories: unique inner fillers. Kan complexes:
all horn fillers. Quasi-categories sit between: existence-only for inner
horns.

\textbf{Homotopy of $1$-simplices.}\quad $f \sim g$ iff there is a
$2$-simplex with edges $f, g, s_0 x$. The relation is an equivalence by
$\Lambda^3$-horn arguments.

\textbf{Homotopy category $\mathrm{ho}(\mathcal{C})$.}\quad Objects
$\mathcal{C}_0$; morphisms $\mathcal{C}_1/\sim$; composition via
inner-horn filler (well-defined modulo $\sim$). $\mathrm{ho}$ is left
adjoint to the nerve.

\textbf{Bridges.}\quad Differential-graded categories (via $\mathrm{N}^{\mathrm{dg}}$)
yield stable quasi-categories; topological / simplicial categories (via
$\mathrm{N}^{\mathrm{coh}}$) yield general quasi-categories. The
$\infty$-category $\mathcal{S}$ of spaces, $\mathcal{S} =
\mathrm{N}^{\mathrm{coh}}(\mathbf{Kan})$, is the universal example.

\textbf{Three models.}\quad Quasi-categories (this book); simplicial
categories; complete Segal spaces. All Quillen-equivalent.
\end{summarybox}


% ═══════════════════════════════════════════════════════════════════════════
\section*{Formal Restatement}
\addcontentsline{toc}{section}{Formal Restatement}
% ═══════════════════════════════════════════════════════════════════════════

\begin{formalrestatement}
\textbf{Quasi-category.}\quad $\mathcal{C} \in \mathbf{sSet}$ with the
property: $\forall n \geq 2, \forall 0 < k < n, \forall \alpha :
\Lambda^n_k \to \mathcal{C}, \exists \tilde\alpha : \Delta^n \to
\mathcal{C}$ extending $\alpha$.

\medskip
\textbf{Horn-type hierarchy.}\quad
\begin{itemize}
  \item Inner horns fill uniquely: $\mathcal{C} = \mathrm{N}(\mathcal{D})$
        for some category $\mathcal{D}$.
  \item Inner horns fill (existence): $\mathcal{C}$ is a quasi-category.
  \item All horns fill: $\mathcal{C}$ is a Kan complex ($\infty$-groupoid).
\end{itemize}

\medskip
\textbf{Homotopy relation.}\quad For $f, g \in \mathcal{C}_1(x, y)$:
$f \sim g \iff \exists \sigma \in \mathcal{C}_2 \text{ with } d_0\sigma = g,
d_1\sigma = f, d_2\sigma = s_0 x$. Equivalence relation.

\medskip
\textbf{Homotopy category.}\quad $\mathrm{ho}(\mathcal{C})$:
$\mathrm{Ob} = \mathcal{C}_0$, $\mathrm{Hom}(x, y) = \mathcal{C}_1(x, y) / \sim$,
$[g] \circ [f] = [d_1 \sigma]$ for any filler $\sigma : \Delta^2 \to \mathcal{C}$
of $(f, g)$, $\mathrm{id}_x = [s_0 x]$.

\medskip
\textbf{Adjunction.}\quad $\mathrm{ho} : \mathbf{QCat}
\rightleftarrows \mathbf{Cat} : \mathrm{N}$, $\mathrm{ho} \dashv \mathrm{N}$.

\medskip
\textbf{$\infty$-category of spaces.}\quad $\mathcal{S} =
\mathrm{N}^{\mathrm{coh}}(\mathbf{Kan})$; $\mathrm{ho}(\mathcal{S}) =
\mathrm{Ho}(\mathbf{Top}) = \mathrm{Ho}(\mathbf{sSet})$.
\end{formalrestatement}
